\section{Introduction}

Portfolio management involves allocating funds among different categories of assets to ensure the proper achievement of a balance between expected return and risk. Modern Portfolio Theory provides the baseline theoretical foundation for this process, arguing that investors should evaluate securities portfolios not solely on their individual merits, but on expected returns, volatility, and pairwise covariance with other assets in the portfolio \autocite{markowitz1952}. Diversification assists investors in reducing portfolio risk without an equivalent reduction in expected return when selected assets exhibit imperfect correlations.

This case study evaluates a diversified portfolio consisting of three Philippine equities: PLDT Inc. (\texttt{TEL}), Manila Electric Company (\texttt{MER}), and Aboitiz Equity Ventures (\texttt{AEV}); two international technology growth equities: NVIDIA Corporation (\texttt{NVDA}) and Meta Platforms, Inc. (\texttt{META}); and one alternative digital asset: Bitcoin (\texttt{BTC}). By combining domestic Philippine equities, Western growth technology stocks, and a cryptocurrency, the asset universe provides a structured framework to evaluate cross-asset, cross-sector, and cross-border diversification benefits. However, individual asset dynamics must be analyzed carefully, as asset volatility varies across market regimes and macro stress periods (Almeida \& Gon\c{c}alves, 2023).

Using daily price data from January 2, 2020 through December 31, 2025, this study examines the risk and return profile of the multi-asset portfolio benchmarked against the SPDR S\&P 500 ETF Trust (\texttt{SPY}). Daily log return series are evaluated to examine descriptive statistics, distributional shapes, tail risk, and correlation dynamics. The analysis compares static Buy-and-Hold execution against systematic Monthly Rebalancing to determine how periodic re-alignment influences portfolio growth, asset weight drift, and downside risk exposure.

\subsection{Risk-Free Rate Benchmark Specification}

The hurdle rate for all excess return, Sharpe ratio, and tangency portfolio calculations is anchored to the Philippine 91-Day Treasury Bill (3-Month T-Bill) secondary market yield. Quoted at an annualized benchmark rate of $R_f = 5.25\%$, it is converted into a daily log-equivalent rate assuming a 252-trading-day annual convention:

\begin{equation}
  r_{f,\text{daily}} = \frac{\ln(1 + R_f)}{252} = \frac{\ln(1 + 0.0525)}{252} \approx 0.00020307 \quad (0.0203\% \text{ per day})
\end{equation}

This hurdle rate represents the riskless baseline return earned by domestic institutional capital.
